% =====================================================================
%  PREVIEW : section "Qu'est-ce que la syntaxe ?" — style conférence
%  Marcos Lahoz Yago — Orange Research, 13 mai 2026
% =====================================================================
\documentclass[aspectratio=169,11pt]{beamer}

% ---------- Theme base ----------
\usetheme{default}
\usecolortheme{default}
\setbeamertemplate{navigation symbols}{}

% ---------- Encoding & fonts ----------
\usepackage[T1]{fontenc}
\usepackage[utf8]{inputenc}
\usepackage{lmodern}
\usepackage{amsmath, amssymb}
\usepackage{booktabs}
\usepackage{tikz}
\usetikzlibrary{positioning, arrows.meta, shapes.geometric, fit, backgrounds, calc}
\usepackage{xcolor}
\usepackage{xspace}

% ---------- Color palette ----------
\definecolor{accent}{RGB}{255, 121, 0}        % Orange
\definecolor{darknavy}{RGB}{26, 43, 76}       % Body text
\definecolor{midgray}{RGB}{120, 130, 140}     % Subtle accents
\definecolor{lightgray}{RGB}{220, 225, 230}   % Backgrounds
\definecolor{frozenblue}{RGB}{124, 163, 208}
\definecolor{probegreen}{RGB}{123, 190, 123}

% ---------- Beamer typography ----------
\setbeamercolor{normal text}{fg=darknavy}
\setbeamercolor{frametitle}{fg=darknavy, bg=}
\setbeamercolor{title}{fg=darknavy}
\setbeamercolor{subtitle}{fg=midgray}
\setbeamercolor{author}{fg=midgray}
\setbeamercolor{date}{fg=midgray}
\setbeamercolor{institute}{fg=midgray}
\setbeamercolor{structure}{fg=accent}
\setbeamercolor{alerted text}{fg=accent}
\setbeamercolor{itemize item}{fg=accent}
\setbeamercolor{itemize subitem}{fg=accent}
\setbeamercolor{block title}{fg=darknavy, bg=lightgray}
\setbeamercolor{block body}{fg=darknavy, bg=lightgray!50}

\setbeamerfont{frametitle}{family=\sffamily, series=\bfseries, size=\large}
\setbeamerfont{title}{family=\sffamily, series=\bfseries, size=\LARGE}
\setbeamerfont{subtitle}{family=\sffamily, series=\mdseries, size=\normalsize}
\setbeamerfont{author}{family=\sffamily, size=\small}
\setbeamerfont{institute}{family=\sffamily, size=\footnotesize}
\setbeamerfont{date}{family=\sffamily, size=\footnotesize}

% ---------- Frametitle: thin rule below ----------
\setbeamertemplate{frametitle}{%
  \nointerlineskip%
  \vspace{0.5em}\hspace{0.6em}\insertframetitle\par%
  \vspace{0.15em}\hspace{0.6em}\textcolor{midgray}{\rule{\dimexpr\paperwidth-1.2em\relax}{0.3pt}}%
  \vspace{0.2em}%
}

% ---------- Footer: name + page ----------
\setbeamertemplate{footline}{%
  \leavevmode%
  \hspace{0.7em}\textcolor{midgray}{\footnotesize Lahoz Yago \textemdash{} Orange Research, 13 mai 2026}\hfill%
  \textcolor{midgray}{\footnotesize \insertframenumber\,/\,\inserttotalframenumber}\hspace{0.7em}%
  \vspace{0.45em}%
}

\setbeamertemplate{itemize item}{\textcolor{accent}{\textbf{\textbullet}}}
\setbeamertemplate{itemize subitem}{\textcolor{accent}{\textendash}}

\newcommand{\mbert}{m\textsc{BERT}\xspace}
\newcommand{\ancora}{UD\_Spanish-AnCora\xspace}

% =====================================================================
\title{La syntaxe espagnole se cache-t-elle dans m\textsc{BERT} ?}
\subtitle{Sondes structurelles et extension isom\'etrique}
\author{Marcos Lahoz Yago}
\institute{Universit\'e Paris-Saclay}
\date{Orange Research \textemdash{} 13 mai 2026}

\begin{document}

% =====================================================================
% Title
% =====================================================================
\begin{frame}[plain]
  \vspace{1.6cm}
  \begin{center}
    {\usebeamerfont{title}\usebeamercolor[fg]{title}\inserttitle}\\[0.4em]
    \textcolor{midgray}{\rule{0.4\paperwidth}{0.3pt}}\\[0.6em]
    {\usebeamerfont{subtitle}\usebeamercolor[fg]{subtitle}\insertsubtitle}\\[1.6em]
    {\usebeamerfont{author}\usebeamercolor[fg]{author}\insertauthor}\\[0.2em]
    {\usebeamerfont{institute}\usebeamercolor[fg]{institute}\insertinstitute}\\[1.4em]
    {\usebeamerfont{date}\usebeamercolor[fg]{date}\insertdate}
  \end{center}
\end{frame}

% =====================================================================
% 1. L'arbre de dépendances
% =====================================================================
\begin{frame}{Le squelette syntaxique : l'arbre de d\'ependances}
  \begin{columns}[T]
    \column{0.52\textwidth}
    \vspace{0.1cm}
    \centering
    \begin{tikzpicture}[scale=0.95, every node/.style={font=\small, color=darknavy}]
      \node (read)    at ( 0,    2.6) {\textbf{read}};
      \node (student) at (-2.0,  1.7) {\textbf{student}};
      \node (book)    at ( 0,    1.7) {\textbf{book}};
      \node (yest)    at ( 2.0,  1.7) {\textbf{yesterday}};
      \node (the1)    at (-2.0,  0.8) {The};
      \node (the2)    at ( 0,    0.8) {the};
      \draw[thick, darknavy] (read) -- (student);
      \draw[thick, darknavy] (read) -- (book);
      \draw[thick, darknavy] (read) -- (yest);
      \draw[thick, darknavy] (student) -- (the1);
      \draw[thick, darknavy] (book)    -- (the2);
      % Edge labels
      \node[midgray, font=\scriptsize] at (-1.15, 2.25) {nsubj};
      \node[midgray, font=\scriptsize] at ( 0.32, 2.25) {obj};
      \node[midgray, font=\scriptsize] at ( 1.15, 2.25) {advmod};
      \node[midgray, font=\scriptsize, anchor=east] at (-2.05, 1.25) {det};
      \node[midgray, font=\scriptsize, anchor=west] at ( 0.05, 1.25) {det};
      % Caption
      \node[midgray, font=\footnotesize, anchor=north] at (0, 0.2)
        {« The student read the book yesterday. »};
    \end{tikzpicture}

    \vspace{0.25cm}
    {\footnotesize\color{midgray}
    \begin{tabular}{@{}l@{\;}c@{\;}l@{}}
      $d_T(\text{read},\text{student})$ & $=$ & $1$\\
      $d_T(\text{The},\text{book})$     & $=$ & $3$\\
      $d_T(\text{The},\text{the})$      & $=$ & $4$\\
    \end{tabular}}

    \column{0.45\textwidth}
    \vspace{0.4cm}
    \begin{itemize}\setlength\itemsep{8pt}
      \item Chaque mot a exactement \alert{une t\^ete} (la racine n'en a aucune).
      \item Les ar\^etes portent des \'etiquettes fonctionnelles
            (\texttt{nsubj}, \texttt{obj}, \ldots) \textemdash{}
            mais \alert{H\&M les ignorent} : seule la topologie compte.
      \item La \alert{distance d'arbre} $d_T(w_i, w_j)$ est le nombre
            d'ar\^etes du chemin unique entre deux mots.
    \end{itemize}

    \vspace{0.4cm}
    {\footnotesize\color{midgray}\itshape
    Cet arbre est l'objet de r\'ef\'erence : on demandera aux
    repr\'esentations de \mbert{} de le reproduire g\'eom\'etriquement.}
  \end{columns}
\end{frame}

% =====================================================================
% 2. "Syntaxe" est un mot piège
% =====================================================================
\begin{frame}{« Syntaxe » : que mesure-t-on exactement ?}
  \vspace{0.2cm}
  Le mot \emph{syntaxe} recouvre plusieurs notions distinctes.
  La sonde structurelle de Hewitt~\& Manning n'en mesure
  \alert{qu'une seule}.

  \vspace{0.5cm}
  \begin{columns}[T]
    \column{0.46\textwidth}
    \begin{block}{\color{accent}Mesur\'e (H\&M)}
      \vspace{0.1cm}
      \begin{itemize}\setlength\itemsep{4pt}
        \item Topologie de l'arbre de d\'ependances UD
              (non \'etiquet\'e, non orient\'e).
      \end{itemize}
      \vspace{1.6cm}
    \end{block}

    \column{0.50\textwidth}
    \begin{block}{\color{midgray}Hors du cadre}
      \vspace{0.1cm}
      \begin{itemize}\setlength\itemsep{2pt}
        \item \'Etiquettes des ar\^etes
              (\texttt{nsubj}, \texttt{obj}, \ldots)
        \item Structure de constituants (NP, VP, PP, \ldots)
        \item Mouvement, liage, contraintes d'\^iles
        \item Accord \`a longue distance (sujet--verbe)
        \item Grammaticalit\'e \textemdash{} ce qui n'est \emph{pas} permis
        \item R\^oles s\'emantiques (agent, patient, \ldots)
      \end{itemize}
    \end{block}
  \end{columns}

  \vspace{0.5cm}
  \small\itshape
  Quand on dit dans ce travail « la syntaxe est encod\'ee dans \mbert{} »,
  on parle uniquement de la \alert{topologie du graphe de d\'ependances}.
  Pas de la syntaxe au sens chomskyen, pas de la grammaticalit\'e,
  pas des relations fonctionnelles.
\end{frame}

% =====================================================================
% 3. L'hypothèse exacte
% =====================================================================
\begin{frame}{L'hypoth\`ese exacte de Hewitt \& Manning (2019)}
  \vspace{0.3cm}
  \centering
  {\Large $\displaystyle \bigl\| B\,(h_i - h_j) \bigr\|_2^{\,2}
   \;\approx\; d_T(w_i,\, w_j)$}

  \vspace{0.7cm}
  \begin{columns}[T]
    \column{0.34\textwidth}
    {\footnotesize
    \textbf{Lecture des symboles}\\[6pt]
    \begin{itemize}\setlength\itemsep{4pt}
      \item $h_i, h_j \in \mathbb{R}^d$ : vecteurs contextuels de
            \mbert{} (couche fix\'ee, $d{=}768$).
      \item $B \in \mathbb{R}^{r \times d}$ : projection
            \alert{lin\'eaire} apprise ($r{=}128 \ll d$).
      \item $d_T$ : distance de graphe dans l'arbre UD annot\'e.
    \end{itemize}}

    \column{0.62\textwidth}
    {\footnotesize
    \textbf{Trois affirmations dans une seule formule}\\[6pt]
    \begin{itemize}\setlength\itemsep{6pt}
      \item \textbf{Lin\'eaire.}
            Si une projection \emph{lin\'eaire} suffit, la syntaxe est
            \alert{d\'ej\`a l\`a} dans la repr\'esentation \textemdash{}
            ce n'est pas la sonde qui l'invente. Une MLP profonde
            la trouverait m\^eme dans des vecteurs al\'eatoires.
      \item \textbf{Sous-espace de basse dimension} ($r \ll d$).
            La syntaxe ne n\'ecessite pas les $768$ axes : elle vit
            sur une petite ardoise dans l'espace \mbert{}.
      \item \textbf{M\'etrique g\'eom\'etrique.}
            Pas seulement « information disponible » : les
            distances euclidiennes au carr\'e \emph{approximent} les
            distances d'arbre. Une structure discr\`ete embarqu\'ee
            comme une m\'etrique continue \textemdash{}
            \alert{c'est le r\'esultat surprenant}.
    \end{itemize}}
  \end{columns}

  \vspace{0.4cm}
  \centering\footnotesize\color{midgray}\itshape
  Mesur\'e par \textbf{Spearman} (corr\'elation des distances) et
  \textbf{UUAS} (ar\^etes correctement reconstruites).
  Hypoth\`ese falsifiable, comparable entre langues et mod\`eles.
\end{frame}

% =====================================================================
% 4. L'arbre vrai est une convention
% =====================================================================
\begin{frame}{L'arbre « vrai » est une convention, pas la v\'erit\'e}
  \vspace{0.2cm}
  L'arbre de r\'ef\'erence est ce qu'une \'equipe d'annotateurs
  (Taul\'e et al., Universitat de Barcelona, AnCora 2008) a marqu\'e en
  suivant les conventions \alert{Universal Dependencies}. Choix sp\'ecifiques :

  \vspace{0.3cm}
  \begin{itemize}\setlength\itemsep{4pt}
    \item Les \emph{contractions} espagnoles
          (\textit{del = de + el}) : repr\'esent\'ees comme deux n\oe{}uds
          sous un rang composite.
    \item Les sujets \emph{pro-drop} : repr\'esent\'es comme des
          n\oe{}uds vides explicites.
    \item La place de la copule (\textit{ser}, \textit{estar}) :
          t\^ete ou d\'ependant ?
    \item L'attachement des modificateurs ambigus (PP-attachment).
  \end{itemize}

  \vspace{0.4cm}
  D'autres formalismes (HPSG, LFG, grammaire de constituants) donneraient
  des arbres \emph{partiellement diff\'erents} pour la m\^eme phrase.

  \vspace{0.4cm}
  \centering
  \alert{Quand on rapporte UUAS $= 0{,}76$, on dit :}\\[2pt]
  \emph{« la sonde concorde \`a $76\,\%$ avec cette convention
  d'annotation »} \textemdash{} pas avec « la syntaxe de
  l'espagnol » dans l'absolu.
\end{frame}

\end{document}
